\subsection{Overview}
The method combines four components: diverse seed construction, exact evaluation of the represented continuous trajectories, two archives with different roles, and a seeded multi-objective variable-neighborhood search. The feasible archive $\mathcal F$ approximates the Pareto front in~\eqref{eq:objectives}. The auxiliary pool $\mathcal P$ retains valid candidates that are outside $\mathcal F$, including communication-infeasible candidates that may be repaired. Every stochastic decision uses one seeded random stream in the coordinator; candidate evaluation may be batched, but integration into $\mathcal F$ and $\mathcal P$ follows generation order.

\begin{figure}[t]
\centering
\fbox{\parbox{0.92\columnwidth}{\centering
Seed construction $\rightarrow$ base selection $\rightarrow$ ruin-and-recreate\\[2pt]
$\downarrow$\\[-2pt]
budget repair $\rightarrow$ continuous evaluation $\rightarrow$ ordered archive update\\[2pt]
$\downarrow$\\[-2pt]
communication repair and neighborhood passes $\rightarrow$ $\mathcal F,\mathcal P$}}
\caption{Method overview. Physical and budget checks precede communication evaluation; evaluated candidates enter the archives in generation order.}
\label{fig:method_overview}
\end{figure}

\subsection{Representation and seed construction}
A candidate stores an ordered tuple of START/TASK/FREE/END occurrences for every agent. Each occurrence contains its position, nonzero boundary velocity, kind, and task identifier when applicable. A physical key excludes occurrence labels, so duplicate route decisions are detected independently of how a move names them. This representation directly supports route-order, assignment, velocity, relay, and free-waypoint edits.

Initialization starts with direct START-to-END routes. Greedy construction repeatedly samples insertions of unvisited tasks, keeps candidates with the largest reward increase, ranks ties by mean completion time, and randomly chooses among the three best ranked candidates. For agents with identical motion limits and initial velocities, cooperative seeds assign the same distance-ordered task subsets to all agents. Partition seeds sort tasks by angle around the start and distribute them cyclically with every agent shift; one additional seed uses a seeded shuffle. Initialization stops after a configured fraction of the time or exact-evaluation budget. These seeds balance immediately communication-feasible convoy routes with geographically separated assignments.

If all possible insertions were enumerated, one greedy construction step would consider $O(NM\bar K)$ positions for average route length $\bar K$. The implementation caps each generated neighborhood stream at $K_N$, so initialization cost is bounded by its configured fraction of the global time and exact-evaluation budgets.

\begin{algorithm}[t]
\caption{TOP-ICC multi-objective variable-neighborhood search}
\label{alg:movns}
\begin{algorithmic}[1]
\Require Problem data, search limits, neighborhood set $\mathcal N$, seed
\Ensure Feasible epsilon archive $\mathcal F$
\State Initialize seeded random stream, $\mathcal F\gets\emptyset$, $\mathcal P\gets\emptyset$
\State Construct and evaluate direct, greedy, cooperative, and partition seeds
\While{time, evaluation, and iteration limits are not exhausted}
    \State Select base $S$ from $ND_{\mathcal F}$ or $ND_{\mathcal P}$
    \State $S'\gets$ \Call{RuinRecreate}{$S$}
    \State Evaluate and, if needed, communication-repair $S'$
    \ForAll{$N\in\mathcal N$}
        \State Apply $N$ to the current passage base and evaluate candidates in order
        \State Repeat while $CV$ decreases, or feasible hypervolume improves, up to $p_{\max}$ passes
    \EndFor
    \State Update the age of the selected pool member and record metrics
\EndWhile
\State \Return $\mathcal F$
\end{algorithmic}
\end{algorithm}

\subsection{Feasible archive, auxiliary pool, and selection}
A physically valid candidate with $CV=0$ is eligible for $\mathcal F$. The archive first removes Pareto-dominated solutions in $(-Q,\bar c)$ and then retains one representative per epsilon box; if capacity is exceeded, both epsilon widths are coarsened. This bounds memory while preserving reward--time coverage.

Candidates not retained in $\mathcal F$, together with solutions displaced from $\mathcal F$, are considered for $\mathcal P$. Pool dominance uses
\begin{equation}
 g(S)=\bigl(CV(S),-Q(S),\bar c(S)\bigr),
 \label{eq:pool_vector}
\end{equation}
so partial communication repair remains visible. The pool removes dominated and over-age members and, if necessary, truncates by crowding distance. Let $ND_{\mathcal F}$ and $ND_{\mathcal P}$ be the nondominated subsets of the two archives. The default probability of selecting $\mathcal P$ is $|ND_{\mathcal P}|/(|ND_{\mathcal F}|+|ND_{\mathcal P}|)$; a two-member crowding tournament then chooses a base within the selected archive. Keeping $CV$ in $\mathcal P$ rather than in the reported objective vector preserves useful stepping stones without rewarding excess connectivity in the final front.

With archive capacities $C_{\mathcal F}$ and $C_{\mathcal P}$, one insertion costs $O(C_{\mathcal F}+C_{\mathcal P}^2)$ in the current implementation: feasible-front dominance is linear, whereas pool dominance is quadratic in the bounded pool size. Memory is $O(C_{\mathcal F}+C_{\mathcal P})$ solutions, excluding cached trajectory data.

\subsection{Perturbation, repair, and neighborhoods}
Ruin-and-recreate perturbs the selected base before local improvement. When the perturbed solution has $CV>0$, a bounded repair operator inserts or modifies FREE visits to reduce communication violation. Candidates that exceed a flight-time budget are repaired before full communication evaluation: removable FREE visits are considered first when their removal preserves local contacts, followed by TASK removals ranked by reward lost per second saved.

The neighborhood set contains whole-agent velocity changes, individual visit-velocity changes, velocity-segment changes, task removal, replacement, cross-agent assignment, reordering, insertion, relay construction, free-waypoint insertion near a visit, free-waypoint insertion near a route segment, and unlinking. Each family samples at most $K_N$ candidates. For a communication-infeasible passage, a family repeats only after strict $CV$ reduction; for a feasible passage, it repeats only after normalized hypervolume gain greater than $\tau$. Both cases are capped at $p_{\max}$ passes. This rule spends additional evaluations only when the family makes measurable progress.

For a batch of $b$ candidates, generation and repair remain sequential. With $w>1$ evaluator processes, at most $\min(b,2w)$ jobs are pending and results are integrated in input order. The default $w=1$ creates no subprocess. Parallel evaluation changes wall time rather than the logical archive order under an evaluation-count budget.

\subsection{Continuous-time evaluator}
For fixed endpoint states, the connector produces piecewise constant-acceleration motion. On each overlap of two motion pieces, relative position is quadratic in time, so squared separation minus $R^2$ has degree at most four. Real roots and interval signs yield maximal range-contact intervals without uniform time discretization. Unchanged routes and pair contacts may be reused from a parent solution when the evaluation context and mission horizon are compatible; otherwise they are rebuilt.

Algorithm~\ref{alg:communication} processes all contact endpoints plus $0$ and $T$. It first records left-limit ages, then exchanges componentwise newest timestamps over every link active at the event. Processing all simultaneous links as a connected component permits instantaneous multi-hop relay. The output contains every $P_{nm}$ and hence $CV$ from~\eqref{eq:communication_constraint}.

\begin{algorithm}[t]
\caption{Event-based temporal communication evaluation}
\label{alg:communication}
\begin{algorithmic}[1]
\Require Piecewise trajectories, range $R$, recurrence bound $T_c$
\Ensure Directed peaks $P_{nm}$ and violation $CV$
\State Extract and merge maximal range-contact intervals for every agent pair
\State Sort unique contact endpoints together with $0$ and $T$
\State Initialize $h_{nm}\gets0$ for $n\ne m$
\ForAll{consecutive event times in chronological order}
    \State Advance source timestamps through components active on the open interval
    \State Update $P_{nm}$ using ages immediately before the event
    \State Exchange newest timestamps through all links active at the event
\EndFor
\State $CV\gets\sum_{n\ne m}\max\{0,P_{nm}-T_c\}$
\State \Return $P, CV$
\end{algorithmic}
\end{algorithm}

\noindent\textbf{Proposition 1.} Given exact maximal range-contact intervals for the represented trajectories, Algorithm~\ref{alg:communication} returns the exact directed peaks $P_{nm}$.

\noindent\emph{Proof sketch.} The instantaneous communication graph is constant on every open interval between consecutive contact endpoints. Within such an interval, each held source timestamp is either propagated through a fixed connected component or remains unchanged, so information age is affine with slope zero or one. Its maximum therefore occurs at an interval boundary, including the left limit before an exchange. At an endpoint, processing all active edges as connected components implements every simultaneous zero-delay relay. Induction over the ordered endpoints establishes the newest timestamp held by every receiver after each event; checking the corresponding left limits and the horizon yields every supremum in~\eqref{eq:staleness}.

If pair $(m,n)$ contains $q_m$ and $q_n$ motion pieces, contact extraction examines $O(q_mq_n)$ overlaps. For $C$ resulting intervals, the chronological sweep has at most $2C+2$ event times; component propagation costs $O(M^2)$ per event in the implemented vector-clock representation. The total evaluation cost is therefore $O(\sum_{m<n}q_mq_n+CM^2)$, apart from constant-degree root solving.

\subsection{Design choices and hyperparameters}
The method uses exact event ordering because an aggregated contact graph can report a path that never exists in chronological order. It uses instantaneous in-range exchange to isolate recurrence from link-layer protocol modeling. The finite velocity menu and sampled neighborhoods make the search tractable, while the full continuous evaluator protects feasibility between visits. The main hyperparameters are the random seed, time and exact-evaluation limits, archive and pool capacities, epsilon widths, velocity headings and speed fractions, $K_N$, $p_{\max}$, hypervolume threshold $\tau$, communication-repair rounds, batch size, and evaluator workers. Section~\ref{sec:experiments} reports the values used in each saved run.
